% Generated draft: requires mathematical and source review.
\documentclass[11pt]{article}
\usepackage[margin=25mm]{geometry}
\usepackage{fontspec}
\setmainfont{Latin Modern Roman}
\usepackage{amsmath,amssymb,mathtools,unicode-math}
\setmathfont{Latin Modern Math}
\usepackage{enumitem,needspace,xurl,hyperref}
\setlength{\parindent}{0pt}
\setlength{\parskip}{5pt}
\setlength{\emergencystretch}{3em}
\newcommand{\N}{\mathbb{N}}
\newcommand{\Z}{\mathbb{Z}}
\newcommand{\Q}{\mathbb{Q}}
\newcommand{\R}{\mathbb{R}}
\newcommand{\C}{\mathbb{C}}
\newcommand{\F}{\mathbb{F}}
\newcommand{\A}{\mathbb{A}}
\newcommand{\PP}{\mathbb{P}}
\newcommand{\E}{\mathbb{E}}
\newcommand{\eps}{\varepsilon}
\newcommand{\dd}{\,\mathrm{d}}
\begin{document}
\section*{UnsolvedMath ID: 800664}
\textbf{Source problem code:} catalogue entry
\subsection{Definitions and mathematical statement}
Let \(X\) be a topological space, and put \(\mathbb{N}=\{1,2,3,\ldots\}\). To interpret the notation in the supplied record, use the following standard cover conventions. An open cover of \(X\) is a family of open subsets whose union is \(X\), with \(X\) itself excluded as a member. For \(n\geq 1\), let \(\mathcal{O}_n(X)\) consist of those open covers \(\mathcal{U}\) such that every subset of \(X\) of cardinality at most \(n\) is contained in some member of \(\mathcal{U}\). Let \(\Gamma(X)\) consist of the infinite open covers \(\mathcal{U}\) such that, for every \(x\in X\), the set \(\{U\in\mathcal{U}:x\notin U\}\) is finite.

The selection property \(\mathsf{S}_1(\{\mathcal{O}_n\}_{n\in\mathbb{N}},\Gamma)\) means that for every sequence \((\mathcal{U}_n)_{n\geq1}\) satisfying \(\mathcal{U}_n\in\mathcal{O}_n(X)\), there are choices \(U_n\in\mathcal{U}_n\) such that \(\{U_n:n\geq1\}\in\Gamma(X)\). The record calls a space satisfying this property a strong \(\gamma\)-set.

In the game \(\mathsf{G}_1(\{\mathcal{O}_n\}_{n\in\mathbb{N}},\Gamma)\), the players ONE and TWO play innings indexed by \(n\geq1\). In inning \(n\), ONE chooses \(\mathcal{U}_n\in\mathcal{O}_n(X)\), and TWO then chooses \(U_n\in\mathcal{U}_n\). TWO wins precisely when \(\{U_n:n\geq1\}\in\Gamma(X)\); otherwise ONE wins. A strategy for ONE specifies each legal move from the preceding finite play, and is winning if every infinite play following it is won by ONE. The selected family must be infinite as a family of distinct sets; repetitions do not themselves satisfy this requirement.

\textbf{Recorded question.} For every \(X\) in the intended class of spaces, is it true that
\[
X\text{ is a strong }\gamma\text{-set}
\quad\Longleftrightarrow\quad
\text{ONE has no winning strategy in }
\mathsf{G}_1(\{\mathcal{O}_n\}_{n\in\mathbb{N}},\Gamma)?
\]
The supplied record does not specify the intended class of spaces or supply the definitions of the cover classes and game. Thus the definitions above make the standard interpretation explicit, but the intended domain and source conventions require confirmation; in particular, the record does not establish whether \(X\) ranges over arbitrary topological spaces or a restricted class such as sets of reals.
\subsection{Short English statement}
The question asks whether being a strong \(\gamma\)-set is equivalent to ONE having no winning strategy in the corresponding cover-selection game. The intended class of spaces and the source's precise cover conventions are not supplied.
\subsection{Sources}
\begin{itemize}
\item Supplied local UnsolvedMath registry record: permanent ID \(800664\), exact source problem code \texttt{catalogue entry}; title: Selection Principles Milestone Problems --- Problem 37 ([strongdiags, Problem 5.16]).
\item Source list attributed to Tsaban: \emph{Selection Principles in Mathematics: A Milestone of Open Problems} (2003), Problem 37. Supplied source URL: \url{https://arxiv.org/abs/math/0312182}. The registry lists Tsaban as proposer and 2003 as the proposed year.
\item Supplied attribution/context: \texttt{[strongdiags, Problem 5.16]}. The supplied record does not expand this bibliographic key.
\item Supplied provenance reports access on 2026-10-16 and extraction by \texttt{arxiv\_latex}. This is recorded provenance, not an independent source visit or verification in this draft.
\item The supplied registry labels the problem \texttt{open}; that label is source metadata, not an independently verified assessment of its current mathematical status.
\end{itemize}
\end{document}
